% GENERATED FILE. Edit canonical lessons, then run npm run generate:books.
% canonical-source-hash: fnv1a64:5ea571e820726f67
% canonical-lessons: IT-C27-sesto, IT-C27-settimo, IT-C27-ottavo, IT-C27-nono, IT-C27-decimo

\chapter{Sesto al Decimo}
\label{ch:it-ordinali-al-decimo}
\hlchaptermodality{27}

\begin{chapteropening}
\textbf{By the end of this chapter:} \emph{I can count in order all the way to tenth in Italian, and say how ottavo and nono each break the shape in a different way.}

The tenth ordinal closes the inherited list, and the dieci/decimo split repeats the one cinque and quinto showed five lessons earlier.
\end{chapteropening}

\section[sesto / sesta]{sesto — the sixth, and a letter that disappeared}

\label{lesson:IT-C27-sesto}



\begin{quote}
Recalls open the chapter, at the four measured distances.

\begin{itemize}
\raggedright
  \item \emph{Recall:} say \emph{quinto} — \textbf{R1}, one lesson back, before it decays
  \item \emph{Recall:} say \emph{primo} — \textbf{R2}, the first real retrieval, five to fifteen lessons back
  \item \emph{Recall:} say \emph{la famiglia} — \textbf{R2}, a second word from the same distance
  \item \emph{Recall:} say \emph{il padre e la madre} — \textbf{R3}, durable, twenty to sixty lessons back
  \item \emph{Recall:} say \emph{il pane} — \textbf{R3}, a second word from the same distance
  \item \emph{Recall:} say \emph{prego, here you are} — \textbf{R4}, recognition at distance, eighty lessons back or more
  \item \emph{Recall:} say \emph{Prego? — come again} — \textbf{R4}, a second word from the same distance
\end{itemize}
\end{quote}

\subsection*{You'll want to know: The word}
\begin{quote}
\textbf{sesto} (masculine) · \textbf{sesta} (feminine) — \textbf{sixth}. \textbf{il sesto piano} — the sixth floor. \textbf{la sesta volta} — the sixth time.
\end{quote}

\begin{cousinweb}
Latin had \emph{sex} and \emph{sextus}, both with an \textbf{x} — which is a \textbf{ks} sound written with one letter. Italian does not keep \emph{x} in inherited words. The cluster came apart and simplified, and both words lost it: \emph{sex} became \textbf{sei}, and \emph{sextus} became \textbf{sesto}.

You can see the same disappearance in every Italian word that had one. It is why Italian spelling looks so much plainer than Latin's, and why a Latin word with an \emph{x} is usually recognisable in Italian with an \emph{s} or a doubled consonant in its place.

English kept the \textbf{x} because it borrowed from books rather than inheriting from speech: \textbf{sextet}, \textbf{sextant}.
\end{cousinweb}

\subsection*{Your turn}
Say these aloud:

\begin{itemize}
\raggedright
  \item \emph{sesto, sesta}
  \item the join — \emph{sei $\to$ sesto}, both from an \emph{x} that went
  \item \emph{la sesta volta} — the sixth time
\end{itemize}

\begin{tcolorbox}[breakable,colback=teal!4,colframe=teal!35!black,title={Before you move on}]
What letter did Latin \emph{sextus} have that \emph{sesto} has not? (\textbf{X} — a \emph{ks} sound.) What happened to the same letter in \emph{sex}? (\textbf{It went too} — \emph{sei}.) And say \textquotedblleft{}the sixth floor.\textquotedblright{} (\textbf{Il sesto piano}.)
\end{tcolorbox}
\section[settimo / settima]{settimo — the seventh, and the seventh heaven}

\label{lesson:IT-C27-settimo}



\begin{quote}
Recalls first, then the seventh.

\begin{itemize}
\raggedright
  \item \emph{Recall:} say \emph{sesto} — \textbf{R1}, one lesson back, before it decays
  \item \emph{Recall:} say \emph{secondo} — \textbf{R2}, the first real retrieval, five to fifteen lessons back
  \item \emph{Recall:} say \emph{il nome} — \textbf{R2}, a second word from the same distance
  \item \emph{Recall:} say \emph{undici a sedici} — \textbf{R3}, durable, twenty to sixty lessons back
  \item \emph{Recall:} say \emph{diciassette a venti} — \textbf{R3}, a second word from the same distance
  \item \emph{Recall:} say \emph{come} — \textbf{R4}, recognition at distance, eighty lessons back or more
  \item \emph{Recall:} say \emph{stare} — \textbf{R4}, a second word from the same distance
\end{itemize}
\end{quote}

\subsection*{You'll want to know: The word}
\begin{quote}
\textbf{settimo} (masculine) · \textbf{settima} (feminine) — \textbf{seventh}. \textbf{il settimo piano} — the seventh floor. \textbf{la settima volta} — the seventh time.
\end{quote}

The double \textbf{tt} stays and you must hold it — \emph{SET-ti-mo}, with the stress right at the front. A single \emph{t} would make a different word's shape.

\subsection*{You'll want to know: al settimo cielo}
\begin{quote}
\textbf{essere al settimo cielo} — to be in seventh heaven, delighted.
\end{quote}

The phrase is old and it is astronomy. Mediaeval writers counted the heavens outward from the earth — Moon, Mercury, Venus, Sun, Mars, Jupiter, and then \textbf{Saturn, the seventh} — and above them the fixed stars. The seventh was the highest a planet reached, so being at it meant being as high as you could get.

English has the same idiom and the same count. Italian says it with the ordinal you have this minute learned, which is a reasonable place to notice that these words are not only for floors.

\subsection*{Your turn}
Say these aloud:

\begin{itemize}
\raggedright
  \item \emph{settimo, settima} — \emph{SET-ti-mo}, hold the \emph{tt}
  \item the join — \emph{sette $\to$ settimo}
  \item \emph{essere al settimo cielo}
\end{itemize}

\begin{tcolorbox}[breakable,colback=teal!4,colframe=teal!35!black,title={Before you move on}]
Where is the stress in \emph{settimo}? (\textbf{On the first syllable}.) What is the seventh heaven, counted out? (\textbf{Saturn} — the outermost planet in the mediaeval count.) And say \textquotedblleft{}the seventh time.\textquotedblright{} (\textbf{La settima volta}.)
\end{tcolorbox}
\section[ottavo / ottava]{ottavo — the ending that breaks step}

\label{lesson:IT-C27-ottavo}



\begin{quote}
Recalls first, then the eighth — the one that does not do what the last four did.

\begin{itemize}
\raggedright
  \item \emph{Recall:} say \emph{settimo} — \textbf{R1}, one lesson back, before it decays
  \item \emph{Recall:} say \emph{terzo} — \textbf{R2}, the first real retrieval, five to fifteen lessons back
  \item \emph{Recall:} say \emph{la persona} — \textbf{R2}, a second word from the same distance
  \item \emph{Recall:} say \emph{nero e bianco} — \textbf{R3}, durable, twenty to sixty lessons back
  \item \emph{Recall:} say \emph{rosso e blu} — \textbf{R3}, a second word from the same distance
  \item \emph{Recall:} say \emph{Come stai?} — \textbf{R4}, recognition at distance, eighty lessons back or more
  \item \emph{Recall:} say \emph{Come sta?} — \textbf{R4}, a second word from the same distance
\end{itemize}
\end{quote}

\subsection*{You'll want to know: The word}
\begin{quote}
\textbf{ottavo} (masculine) · \textbf{ottava} (feminine) — \textbf{eighth}. \textbf{l'ottavo piano} — the eighth floor. \textbf{l'ottava volta} — the eighth time.
\end{quote}

\emph{Otto} is eight. By now you would expect \emph{otto} with the ending the last four took, and it is not that: this one ends in \textbf{-avo}, and it is the only one of the ten that does.

\begin{cousinweb}
Latin \emph{octāvus} already had the odd \textbf{-āvus}, and Italian inherited the oddity rather than tidying it away — which is what usually happens to a word that is used too often to be regularised.

\textbf{L'ottava} in music is the \textbf{eighth} note: \emph{do re mi fa sol la si} and then \emph{do} again, where the names come round. The interval is named for the count that closes it, and Italian named the notes too — \emph{do, re, mi} are the opening syllables of a mediaeval hymn, set down by an Italian monk.

So the word for a musical octave and the word for the eighth floor are the same word, and the music is the one that explains the name.
\end{cousinweb}

\subsection*{Your turn}
Say these aloud:

\begin{itemize}
\raggedright
  \item \emph{ottavo, ottava}
  \item the odd ending — \textbf{-avo}, not the ending the last four took
  \item \emph{l'ottava} — the eighth note, where the names begin again
\end{itemize}

\begin{tcolorbox}[breakable,colback=teal!4,colframe=teal!35!black,title={Before you move on}]
What ending does \emph{ottavo} take, and how many of the ten take it? (\textbf{-avo}, and \textbf{only this one}.) What is an octave, counted out? (\textbf{The eighth note}, where the names start again.) And say \textquotedblleft{}the eighth floor.\textquotedblright{} (\textbf{L'ottavo piano}.)
\end{tcolorbox}
\section[nono / nona]{nono — the ninth, squeezed shut}

\label{lesson:IT-C27-nono}



\begin{quote}
Recalls first, then the ninth — the one that gets shorter.

\begin{itemize}
\raggedright
  \item \emph{Recall:} say \emph{ottavo} — \textbf{R1}, one lesson back, before it decays
  \item \emph{Recall:} say \emph{quarto} — \textbf{R2}, the first real retrieval, five to fifteen lessons back
  \item \emph{Recall:} say \emph{il cuore} — \textbf{R2}, a second word from the same distance
  \item \emph{Recall:} say \emph{avere} — \textbf{R3}, durable, twenty to sixty lessons back
  \item \emph{Recall:} say \emph{ho parlato} — \textbf{R3}, a second word from the same distance
  \item \emph{Recall:} say \emph{Come va?} — \textbf{R4}, recognition at distance, eighty lessons back or more
  \item \emph{Recall:} say \emph{così così} — \textbf{R4}, a second word from the same distance
\end{itemize}
\end{quote}

\subsection*{You'll want to know: The word}
\begin{quote}
\textbf{nono} (masculine) · \textbf{nona} (feminine) — \textbf{ninth}. \textbf{il nono piano} — the ninth floor. \textbf{la nona volta} — the ninth time.
\end{quote}

\begin{cousinweb}
\emph{Nove} is nine. Every ordinal so far has been at least as long as its cardinal; this one is \textbf{shorter}. Latin had \emph{novēnus} and squeezed it shut into \emph{nōnus}, and Italian inherited the squeezed form.

Two ordinals running have now bent the shape in two different directions — \emph{ottavo} by taking an odd ending, \emph{nono} by losing sound. That is what a series of inherited words looks like from inside: not a rule with exceptions, but ten words that each survived a slightly different accident.

And the feminine stands on its own in Italian. \textbf{La Nona} — no noun after it — means Beethoven's ninth symphony, the way English says \emph{the Ninth}.
\end{cousinweb}

\subsection*{Your turn}
Say these aloud:

\begin{itemize}
\raggedright
  \item \emph{nono, nona}
  \item the squeeze — \emph{nove $\to$ nono}, shorter than its number
  \item \emph{la Nona} — the Ninth
\end{itemize}

\begin{tcolorbox}[breakable,colback=teal!4,colframe=teal!35!black,title={Before you move on}]
Is \emph{nono} longer or shorter than \emph{nove}? (\textbf{Shorter} — it was squeezed shut.) Which two ordinals bend the shape, and how differently? (\emph{\textbf{Ottavo}} takes an odd ending; \emph{\textbf{nono}} loses sound.) And what is \emph{la Nona}? (\textbf{The Ninth} — the symphony.)
\end{tcolorbox}
\section[decimo / decima]{decimo — the tenth, and the split heard again}

\label{lesson:IT-C27-decimo}



\begin{quote}
Recalls first, then the tenth, which closes the memorized ten.

\begin{itemize}
\raggedright
  \item \emph{Recall:} say \emph{nono} — \textbf{R1}, one lesson back, before it decays
  \item \emph{Recall:} say \emph{quinto} — \textbf{R2}, the first real retrieval, five to fifteen lessons back
  \item \emph{Recall:} say \emph{l'occhio} — \textbf{R2}, a second word from the same distance
  \item \emph{Recall:} say \emph{parlò} — \textbf{R3}, durable, twenty to sixty lessons back
  \item \emph{Recall:} say \emph{essere} — \textbf{R3}, a second word from the same distance
  \item \emph{Recall:} say \emph{non c'è male} — \textbf{R4}, recognition at distance, eighty lessons back or more
  \item \emph{Recall:} say \emph{the second chapter's own recombining} — \textbf{R4}, a second word from the same distance
\end{itemize}
\end{quote}

\subsection*{You'll want to know: The word}
\begin{quote}
\textbf{decimo} (masculine) · \textbf{decima} (feminine) — \textbf{tenth}. \textbf{il decimo piano} — the tenth floor. \textbf{la decima volta} — the tenth time.
\end{quote}

Stress at the front: \emph{DE-chi-mo}, three syllables.

\begin{grammarlens}[title={What you've built: the split, and the end of the list}]
The \emph{quinto} lesson promised this pair, and here it is. \textbf{Dieci} has a soft \textbf{ch}-sound in the middle; \textbf{decimo} has the same letter and it is soft too — \emph{DE-chi-mo} — but the vowel went the other way: the cardinal broke its \emph{e} into \textbf{ie} and the ordinal did not. One more word that moved while its twin stood still.

\textbf{La decima} was the \textbf{tithe}, the tenth of a harvest owed to the church, and Italian still uses the feminine that way.

That is the ten. \emph{Primo} to \emph{decimo} — ten separate words, each inherited whole from Latin, each with its own small accident. There is no rule that generates them, and you have now met all of them. The next lesson is the good news: from eleven upward Italian stopped inheriting and started building, and one ending does the rest of the number line.
\end{grammarlens}

\subsection*{Your turn}
Say these aloud:

\begin{itemize}
\raggedright
  \item \emph{decimo, decima} — \emph{DE-chi-mo}
  \item the pair — \emph{dieci} broke its vowel, \emph{decimo} did not
  \item all ten in order, \emph{primo} to \emph{decimo}
\end{itemize}

\begin{tcolorbox}[breakable,colback=teal!4,colframe=teal!35!black,title={Before you move on}]
How is \emph{decimo} said? (\emph{\textbf{DE-chi-mo}}, stress at the front.) What did \emph{dieci} do to its vowel that \emph{decimo} did not? (\textbf{Broke it into \emph{ie}}.) And how many separate inherited words are there from first to tenth? (\textbf{Ten} — no rule generates them.)
\end{tcolorbox}
